\documentclass[11pt]{article}
\usepackage[T1]{fontenc}
\usepackage{lmodern,amsmath,amssymb,amsthm,mathtools}
\usepackage[margin=1in]{geometry}
\usepackage[colorlinks=true,linkcolor=blue,citecolor=blue,urlcolor=blue]{hyperref}
\usepackage{microtype}
\newtheorem{theorem}{Theorem}[section]
\newtheorem{lemma}[theorem]{Lemma}
\newtheorem{corollary}[theorem]{Corollary}
\theoremstyle{remark}\newtheorem{remark}[theorem]{Remark}
\newcommand{\osc}{\operatorname{osc}}
\newcommand{\tr}{\operatorname{tr}}
\newcommand{\R}{\mathbb R}
\newcommand{\FIT}{\mathrm{FIT}}
\title{Uniform interior gradient estimates for finite $p$\texorpdfstring{\\}{ }through a constrained affine lift}
\author{Myeongjin Kang\thanks{Copyright 2026 Myeongjin Kang. Licensed under \href{https://creativecommons.org/licenses/by/4.0/}{CC BY 4.0}.}}
\date{October 14, 2026}
\hypersetup{pdfauthor={Myeongjin Kang},pdftitle={Uniform interior gradient estimates for finite p through a constrained affine lift}}
\begin{document}
\maketitle

\begin{abstract}
Assuming the interior regularity and fitted-class Liouville results
of the pinned OpenAI infinity-harmonic manuscript, we give an argument
for a dimension-dependent interior $C^{1,\alpha}$ estimate common to
all finite $p\geq3$. The argument uses the exact equation for the
affine lift $u(y)+b$, including at critical points of $u$. Its
anisotropic limits have an effective viscosity. Finite viscosity is
handled by fitted-class rigidity, whose positive-viscosity proof is
included. Infinite viscosity is handled by a uniform small-solution
estimate and a quadratic comparison between different fitting scales.
The exponent is nonexplicit. Supporting Lean files certify the
zero-jet vector ellipticity algebra, not the analytic theorem.
\end{abstract}

\section{Statement and inputs}

For a symmetric matrix $X$ and vector $v$, write $X[v]=v^TXv$.
An operator with positive Hessian sign is nonnegative at upper tests
of a viscosity subsolution and nonpositive at lower tests of a
supersolution. Weak $p$-harmonic functions are identified with their
continuous representatives.

\begin{theorem}\label{thm:main}
For every integer $n\geq2$ there are $\alpha_n\in(0,1/3]$ and
$C_n<\infty$ such that for every finite $p\geq3$ and every bounded
weak solution $u\in W^{1,p}_{\mathrm{loc}}(B_1\subset\R^n)$ of
\[
 \int |Du|^{p-2}Du\cdot D\phi=0\qquad(\phi\in C_c^\infty(B_1)),
\]
one has $u\in C^{1,\alpha_n}(B_{1/2})$ and
\begin{equation}\label{eq:main}
 \|Du\|_{L^\infty(B_{1/2})}+[Du]_{C^{0,\alpha_n}(B_{1/2})}
 \leq C_n\osc_{B_1}u.
\end{equation}
There is no hypothesis on sign, critical points, level sets or boundary data.
\end{theorem}

The theorem is relative to the stated results of \cite{OpenAI} at commit
\begin{center}\texttt{fd4aeeb2ee4fc729c18d98444fed42fd0529eeeb}.\end{center}
Specifically,
we use its infinity-harmonic interior estimate in dimensions at least
three, its exact infinity-harmonic linear lift, its zero-viscosity
fitted-class Liouville theorem, and its equation-free geometric fit
and envelope lemmas. The new PDE identifications are proved here.
We do not assume that an affine lift of a finite-$p$ solution is
$p$-harmonic. The inhomogeneous equation and the range $2<p<3$ are
outside the statement.

The classical analytic inputs are the local logarithmic gradient
estimate of Kotschwar--Ni \cite{KotschwarNi}, scalar regularity from
Uhlenbeck \cite{Uhlenbeck}, weak/viscosity equivalence
\cite{JulinJuutinen}, Savin's small-solution theorem \cite{Savin}, and
linear weak-solution Harnack \cite{Crescimbeni}. The latter is used
in Appendix~\ref{app:viscous}. Their conclusions are not new claims
of this manuscript. A companion audit records exact source versions,
read scopes and the compact-parameter choices in Uhlenbeck's proof.

Appendix~\ref{app:viscous} supplies the positive-viscosity argument:
the cone infimum, viscosity-to-weak passage, contact propagation and
rigidity. No estimate uniform in the positive viscosity is used.

\paragraph{Relation to existing estimates.}
Peng--Zhang--Zhou \cite[Theorem 1.1]{PengZhangZhou} establish
$C^1$ convergence for potentials on convex rings, with geometric
and nonvanishing-gradient structure. Bungert \cite[Theorem 1]{Bungert}
gives value-convergence rates under a value-Hölder hypothesis.
Neither statement supplies a common gradient Hölder estimate for
arbitrary interior weak solutions. The qualitative-to-quantitative
theorem of Sobral--Teixeira--Urbano
\cite[Theorem 1.1]{SobralTeixeiraUrbano} gives constants depending on
$p$. Its large-shift compactness at fixed $p$ produces a linear
operator with matrix $I+(p-2)e\otimes e$, whose ellipticity ratio
diverges in the joint large-$p$ limit. The effective-viscosity
analysis below addresses this simultaneous degeneracy.

Lee--Lee \cite{LeeLee} prove improved homogeneous regularity and the
inhomogeneous $C^{p'}$ conclusion for $p$ sufficiently close to two.
Their Theorem A transfers a homogeneous exponent strictly larger than
$1/(p-1)$ to the sharp bounded-forcing exponent. If
Theorem~\ref{thm:main} is established, that transfer applies for each
fixed $p>\max\{3,1+1/\alpha_n\}$. It is a corollary of the present
claim, with constants allowed to depend on $p$; the intervening range
is not settled. In particular the full inhomogeneous conjecture is
not claimed here.

\section{Uniform Lipschitz bounds and qualitative flatness}

For a positive $p$-harmonic function $v$ on a Euclidean ball $B_R$,
put $w=-(p-1)\log v$. In \cite[Theorem 1.1]{KotschwarNi}, choose
the auxiliary parameter $1/2$ and zero curvature bound. For $p\geq3$
the displayed constants reduce to
\[
 \begin{split}
 c_{p,n}&=(p/2+1)(p-1),\qquad
 b_{p,n}=2(p-1)/(n-1)-2,\\
 C(n,0,p,R,1/2)&=
 \frac{20(n+p-2)(n-1)+10(p-1)(n-1)}{R^2}.
 \end{split}
\]
Substitution gives $|Dw|\leq K_n(p-1)/R$ on $B_{R/2}$, hence
$|D\log v|\leq K_n/R$. If $M=\osc u>0$, replace $u$ by
$v=u-\inf u+M$, so $M\leq v\leq2M$ and
$|Du|\leq2K_nM/R$. The case $M=0$ is immediate. Classical scalar
regularity gives the differentiable representative; the gradient
estimate's smooth noncritical argument extends across critical points
where the gradient is zero.

Normalize $\osc_{B_1}u\leq1$ and $u(0)=0$. The affine lifts
\[
 U(y,b)=u(y)+b,\qquad \xi=e_b,
\]
are uniformly bounded and $L$-Lipschitz on $B_{7/8}\subset\R^{n+1}$,
with $L$ depending only on $n$, and satisfy
$U(z+h\xi)=U(z)+h$.

\begin{lemma}[Compact-$p$ interface]\label{lem:compact-p}
For every $n\geq2$ and finite $P\geq3$, there are common
$\theta(n,P)>0$ and $C(n,P)<\infty$ such that every bounded weak
$p$-harmonic solution on $B_1$, $3\leq p\leq P$, satisfies
\[
 \|Du\|_{L^\infty(B_{1/2})}+[Du]_{C^{0,\theta}(B_{1/2})}
 \leq C(n,P)\osc_{B_1}u.
\]
This includes zero gradients. An entire weak solution for one fixed
$p\geq3$ with bounded gradient is affine.
\end{lemma}
\begin{proof}
We give the common choices in \cite[Sections 1--5]{Uhlenbeck}; a
fixed-$p$ estimate alone would not establish uniformity. Work first
in dimension $m\geq3$, and distinguish Uhlenbeck's growth exponent
$s=(p-2)/2\in[1/2,S]$, $S=(P-2)/2$, from the present $p$. Set
$\omega=Du$, $Q=|\omega|^2$, $\rho(Q)=Q^s$, and set her structural
constant $C$ to zero. The gradient/curl complex satisfies her (1.2),
and $\omega$ solves her system (1.1).

\paragraph{Common structural and preliminary estimates.}
Fix
\[
 a=\min\{1/2,1/(m-2)\},\qquad
 K_*=2+(1+2S)+S\max\{1,S\}.
\]
The exponent $a$ is positive, at most every $s$, and strictly below
$2/(m-2)$, so $2(1+a)$ is one fixed admissible Sobolev exponent.
The ellipticity expression is
$\rho+2Q\rho'=(1+2s)Q^s$, and
\begin{equation}\label{eq:compact-structural}
 |s x^s-s y^s|\leq S\max\{1,S\}(x+y)^{s-a}|x-y|^a
 \qquad(x,y\geq0).
\end{equation}
Indeed divide by $(x+y)^s$ when $x+y>0$. On $[0,1]$, the power
$t^s$ is $s$-H\"older with constant one for $s\leq1$, and is
$S$-Lipschitz for $s\geq1$. Since $a\leq s$ and $a\leq1$, either
bound implies \eqref{eq:compact-structural}; the zero case is direct.
Thus (1.3), (1.3)$'$ and (1.4) of the source hold with $K_*$.
Differentiability of $\rho$ is required only at $Q>0$.

The source's (1.7), (1.10) and Sections 2--3 apply to
\[
 H_s(Q)=\frac{s+1/2}{s+1}Q^{s+1}
       =\frac{p-1}{p}|Du|^p.
\]
Its prefactor is in $[2/3,1]$. The divergence-form operator in
(1.10) has common measurable ellipticity bounds controlled by
$K_*$. Fix a mean-value exponent $\beta$ strictly between $1$ and
$m/(m-2)$ in (2.1), and fix the Sobolev exponent. Moreover,
\[
 |D\sqrt{H_s(Q)}|\leq(S+1)Q^{s/2}|D\omega|.
\]
The weighted differentiability and energy estimate (3.1), followed
by these fixed Sobolev and Moser estimates in the proof of (3.2),
give one $C_2(m,P)$ such that
\begin{equation}\label{eq:compact-moser}
 \sup_{B_r}H_s(Q)\leq C_2(m,P)r^{-m}
                         \int_{B_{4r}}Q^{s+1}.
\end{equation}
The constants here and in (3.3) involve $K_*$, $(S+1)^2$, fixed
cutoff factors, Sobolev and Poincar\'e constants, and Moser's
$c(\beta)$ and weak-Harnack constant $c'(1)$. None requires a
parameter-dependent threshold choice. In the scalar case the
potential for $\omega$ is $u$ itself; no varying-exponent Hodge
potential estimate is implicit. The weak differentiability and
subsolution justification in Sections 2--3 includes $Q=0$.

\paragraph{One common small-excess iteration.}
Fix the relative reference size $\eta=1/2$ in the source's (4.1).
Normalize $\sup Q=M=1$ by homogeneity; $M=0$ is trivial. During the
iteration maintain $|w_i|\leq1$ and $|w_i|^2\geq1/4$. The
normalized constant-coefficient linearization at a reference $w$ is
\[
 I+2s\frac{w}{|w|}\otimes\frac{w}{|w|},
\]
with eigenvalues in $[1,P-1]$. Its comparison energy and interior
estimates in (4.2) therefore have common constants.

Here is a uniform derivation of the nonlinear error estimate,
without leaving the power-map constant of printed (4.6) implicit.
Put
\[
 A_s(v)=|v|^{2s}v,\qquad F_s(v)=|v|^s v,\qquad
 G_s(v,w)=\frac{A_s(v)-A_s(w)-DA_s(w)(v-w)}{\rho(|w|^2)}.
\]
For $q=2s\in[1,2S]$, the second derivative of $A_s$ away from zero
is bounded by $C_mq(q+3)|v|^{q-1}$. Integration on a segment,
split at zero when necessary, shows that $DA_s$ is uniformly
Lipschitz on the unit ball. Since $\rho(|w|^2)\geq4^{-S}$,
\begin{equation}\label{eq:compact-remainder}
 |G_s(v,w)|\leq C(m,P)|v-w|^{1+a}
 \quad (|v|,|w|\leq1,\ |w|^2\geq1/4).
\end{equation}
To invert the power map uniformly, let $R=\max\{|v|,|w|\}$.
The inequality $DF_s(z)\geq|z|^sI$, integrated over the quarter of
the joining segment nearest a maximal endpoint, gives
\[
 (F_s(v)-F_s(w))\cdot(v-w)
       \geq2^{-S-2}R^s|v-w|^2.
\]
Here that segment has norm at least $R/2$, while $R\geq1/2$.
Cauchy--Schwarz consequently yields
\begin{equation}\label{eq:compact-inverse}
 |v-w|\leq2^{2S+2}|F_s(v)-F_s(w)|.
\end{equation}

Let $E=\int|\omega-w|^2$ on a fixed outer ball. The weighted
estimate (4.8) has a common scalar constant: $DA_s(\omega)$ has
eigenvalues between $Q^s$ and $(P-1)Q^s$, so testing the
differentiated equations against cutoff squared times
$\omega_j-w_j$ gives a weighted-energy factor at most
$4(P-1)^2$ times fixed cutoff constants. The source's
difference-quotient argument justifies this at weak solutions and
zeros of $\omega$. Also
\[
 \|DF_s(v)\|\leq(S+1)|v|^s,\qquad
 |F_s(v)-F_s(w)|\leq(S+1)|v-w|
 \quad (|v|,|w|\leq1).
\]
Thus on nested fixed balls,
\[
 \|F_s(\omega)-F_s(w)\|_{W^{1,2}}^2\leq C(m,P)E.
\]
The fixed embedding into $L^{2(1+a)}$, together with
\eqref{eq:compact-remainder}--\eqref{eq:compact-inverse} and the
linear comparison tested against its zero-boundary difference,
shows that the squared $L^2$ error of the linearized comparison
field is at most $C(m,P)E^{1+a}$, as required in (4.3)--(4.4).

Specify the next reference as follows. If $\widetilde\omega_i$ is
the linearized comparison to $\omega^i-w_i$ on the fixed inner
ball, take
\[
 \widehat w_{i+1}=w_i+\widetilde\omega_i(0),\qquad
 w_{i+1}=\operatorname{proj}_{\overline B_1}(\widehat w_{i+1}).
\]
The center value is controlled by the common interior estimate;
its gradient bound gives the $O(t^2)$ excess improvement. Projection
does not increase the error to any $\omega^i(x)\in\overline B_1$
or the distance to $w_i\in\overline B_1$. With the fixed half-ball
rescalings in (4.9)--(4.11) absorbed into one $C_3=C_3(m,P)\geq1$,
let $E_i$ be the $L^2$ excess on the rescaled unit ball and
$\omega^i$ the corresponding rescaled field. We obtain
\begin{equation}\label{eq:compact-iteration}
 E_{i+1}\leq C_3(t^2+t^{-m}E_i^a)E_i,\qquad
 |w_{i+1}-w_i|\leq C_3\sqrt{E_i}.
\end{equation}
Choose once
\begin{equation}\label{eq:compact-choices}
 \begin{split}
 t&=\min\{1/8,1/(4C_3)\},\\
 \epsilon&=\min\left\{1,\frac1{256C_3^2},
              \left(\frac{t^{m+1}}{4C_3}\right)^{1/a}\right\}.
 \end{split}
\end{equation}
If $E_0\leq\epsilon$, then $E_{i+1}\leq(t/2)E_i\leq tE_i$,
and
\[
 \sum_i|w_{i+1}-w_i|
 \leq\frac{C_3\sqrt\epsilon}{1-\sqrt t}
 \leq2C_3\sqrt\epsilon\leq1/8.
\]
All references have norm at most one, so their squared norms differ
from $|w_0|^2\geq1/2$ by at most twice this drift. Hence
$|w_i|^2\geq1/4$, closing the induction. The source's
(4.12)--(4.15), with their common geometric-series ratio, now give
(4.1) with one threshold $\epsilon_*>0$ and one $1/2$-H\"older
constant $K_4(m,P)$. Fixed unit-ball versus $B_2/B_1$ conventions
only decrease $\epsilon_*$ by fixed radius and volume factors.

\paragraph{One common decrease factor, including critical points.}
Proposition (5.1) uses (3.3), the fixed weak-Harnack constant
$c'(1)$, Poincar\'e for $F_s(\omega)$, and the factors
$(s+1)^2$ and $2^{2s+2}$. These are bounded by $(S+1)^2$ and
$2^{2S+2}$; all other factors are the common $K_*$, cutoff,
Sobolev and Moser constants just listed. The radial maps selecting
the reference have exponent $s+1\in[3/2,S+1]$, so neither their
exponents nor their reciprocals introduce a degenerating choice.
Choose one $C_5(m,P)\geq1$ for that proposition. The two required
inequalities on printed p.~238 are
\[
 1-C_5\sqrt\lambda\geq2^{-(s+1)/2},\qquad
 \lambda C_5 2^{s+1}\leq\epsilon_*.
\]
They hold simultaneously for all $s$ with
\[
 \lambda_*=\min\left\{1/4,\frac1{16C_5^2},
                       \frac{\epsilon_*}{2C_5 2^{S+1}}\right\}>0.
\]
Indeed the first left side is at least $3/4$, while its right side
is at most $2^{-3/4}<3/4$; the second inequality is direct. Set
\[
 \theta=-\frac{\log(1-\lambda_*)}{2\log4}>0,
 \qquad\sqrt{1-\lambda_*}=4^{-\theta}.
\]
Then $\theta<1/2$. Repeated maximum decrease in (5.4)--(5.5)
gives the $r^\theta$ modulus; at the first scale where the
perturbation alternative applies, the common $1/2$-H\"older bound
is stronger. If the decrease continues forever, the maximum tends
to zero, covering a critical gradient without linearizing at zero.
Consequently
\begin{equation}\label{eq:compact-local-holder}
 [\omega]_{C^{0,\theta}(B_{r/2})}
       \leq K_0(m,P)r^{-\theta}\sqrt M
 \quad\text{if }\sup_{B_r}|\omega|^2\leq M.
\end{equation}
Translation supplies the estimate at every center.

\paragraph{The oscillation bound and dimension two.}
Normalize $0\leq u\leq1$ by subtracting its essential infimum and
dividing by its essential oscillation; zero oscillation is trivial.
A cutoff $\chi=1$ on $B_{7/8}$, supported in $B_1$, with
$|D\chi|\leq16$, is used in the admissible weak test $\chi^pu$.
H\"older's inequality, with the zero-energy case separated, gives
\[
 \int\chi^p|Du|^p\leq
 p\left(\int\chi^p|Du|^p\right)^{(p-1)/p}
   \left(\int|u|^p|D\chi|^p\right)^{1/p}.
\]
Sobolev approximation justifies the test. Therefore
\[
 \int_{B_{7/8}}|Du|^p\leq(16P)^P|B_1|=:E_*.
\]
Apply \eqref{eq:compact-moser} with centers in $B_{3/4}$ and
$r=1/32$, so $B_{4r}$ stays in $B_{7/8}$. It gives
\[
 \sup_{B_{3/4}}|Du|^p\leq(3/2)C_2 32^m E_*=:A_*,\qquad
 \|Du\|_{L^\infty(B_{3/4})}\leq\max\{1,A_*\}^{1/3}=:L_*.
\]
Use \eqref{eq:compact-local-holder} on $B(x,1/4)$ for
$x\in B_{1/2}$. Pairs separated by at most $1/8$ obey that local
bound; more distant pairs obey
$|Du(x)-Du(y)|\leq2\,8^\theta L_*|x-y|^\theta$. This proves the
claimed gradient norm and full seminorm after restoring oscillation.
For original dimension two, extend $u$ constantly in a third
coordinate. Fubini shows it is an actual weak $p$-harmonic function
on $B_1^2\times\R$, with the same gradient norm and oscillation.
Apply the three-dimensional estimate on $B_1^3$ and restrict back.
This constant extension is distinct from the affine lift $u+b$.
Finally, for a single fixed $p$, scale
\eqref{eq:compact-local-holder} on arbitrarily large balls for an
entire bounded-gradient solution. Its gradient oscillation on any
fixed ball tends to zero, so it is affine, also as in the source's
Bernstein consequence (6.1).
\end{proof}


The normalized family has a common $o(r)$ Euclidean affine-fit error
at interior points. Otherwise choose $p_i\geq3$, $y_i\in B_{3/4}$,
$r_i\downarrow0$ and rescale
\[
 f_i(Y)=\frac{u_i(y_i+r_iY)-u_i(y_i)}{r_i}
\]
so they remain a fixed positive distance from every linear function
on $B_2$. The common Lipschitz bound gives a locally uniform entire
limit $f$, with $f(0)=0$. After a subsequence $p_i\to p_*<\infty$
or $p_i\to\infty$. Weak/viscosity equivalence identifies the
continuous polynomial equation
\begin{equation}\label{eq:p-poly}
 D^2u[Du]+\epsilon |Du|^2\Delta u=0,
 \qquad\epsilon=(p-2)^{-1}.
\end{equation}
At zero test gradient both this expression and the continuous
$p$-Laplacian vanish. Viscosity stability therefore gives a weak
$p_*$-harmonic limit in the finite case, affine by
Lemma~\ref{lem:compact-p}, or an infinity-harmonic limit otherwise.
In the latter case $F(Y,b)=f(Y)+b$ is infinity-harmonic in dimension
$n+1\geq3$ by the pinned exact lift. Scaling its interior estimate
on $B_R$ gives $[DF]_{\alpha;B_{R/2}}\leq C\operatorname{Lip}(F)R^{-\alpha}$.
Letting $R\to\infty$ makes $F$, hence $f$, affine. Both alternatives
contradict the chosen error. Adding back the exact $b$ coordinate
gives the same qualitative flatness for $U$.

\section{The exact lifted equation}

Put $\Pi=I-\xi\xi^T$. The lift satisfies, at all viscosity tests,
\begin{equation}\label{eq:lift}
 \mathcal L(DU,H)=H[DU]+\epsilon|\Pi DU|^2\tr(\Pi H)=0.
\end{equation}
Here is a proof covering critical points. At an upper test $\phi$
at the origin, restriction to the $b$ line gives $\phi_b=1$ and
$H_{bb}\geq0$. For $\sigma>0$ replace $\phi$ by
$\phi_\sigma=\phi+\sigma b^2/2$. The implicit equation
$(\phi_\sigma)_b=1$ defines a local $C^1$ branch $b=b(y)$.
The function $\psi(y)=\phi_\sigma(y,b(y))-b(y)$ is a $C^2$ upper
test of $u$, with gradient $g=D_y\phi$ and Hessian
\[
 K=H_{yy}-\frac{H_{yb}H_{by}}{H_{bb}+\sigma}.
\]
If $H_\sigma=H+\sigma\xi\xi^T$, then
$H_\sigma-\operatorname{diag}(K,0)$ is positive semidefinite, being
the outer product of $(H_{yb},H_{bb}+\sigma)$ divided by
$H_{bb}+\sigma$. The matrix
$M=(g,1)(g,1)^T+\epsilon|g|^2\Pi$ is positive semidefinite, and
\[
 \tr(MH_\sigma)\geq\tr((gg^T+\epsilon|g|^2I)K)\geq0.
\]
Let $\sigma\downarrow0$; no convergence of $K$ is needed. When
$g=0$, the right-hand side is zero and the vertical curvature is
sufficient. Lower tests use $\phi-\sigma b^2/2$ and the negative
semidefinite block comparison. This proves \eqref{eq:lift}.

\section{Cylinder selection and compactness}

Write $N=n+1$, $G=B_{3/4}\subset\R^N$ and
$D(z)=\operatorname{dist}(z,\partial G)$. For nonzero $P$ put
$N_P=P/|P|$ and
\[
 Z_\lambda(z,r;P)=
 \{z+h:|h\cdot N_P|\leq r,\
 |h-(h\cdot N_P)N_P|\leq\sqrt\lambda r\}.
\]
A parameter $\lambda\in(0,1)$ is good at $(U,z,r)$ if some nonzero
$P$ satisfies
\begin{equation}\label{eq:fit-physical}
 |U(z+h)-U(z)-P\cdot h|\leq\lambda r
 \qquad(z+h\in Z_\lambda(z,r;P)).
\end{equation}
Let $\beta(U,z,r)$ be the supremum of bad parameters, with supremum
zero for the empty set, and let $b(s)$ be the supremum of these
quantities over the normalized family, $z\in G$, and
$0<r\leq sD(z)$. We only use $s\leq1/8$, so all cylinders lie in
the solution domain. The inscribed ball and exact affine increment
imply, for every fit,
\begin{equation}\label{eq:slope-bounds}
 |P|\leq L+\sqrt\lambda,\qquad |P\cdot\xi-1|\leq\sqrt\lambda.
\end{equation}
Thus sufficiently accurate fits have $1/2\leq|P|\leq L+1$.
The qualitative flatness above implies $b(s)\to0$: otherwise a bad
parameter bounded below persists at radii tending to zero, whereas
an affine limit with added-coordinate slope one makes that parameter
good. This argument does not require monotonicity of goodness.

\begin{lemma}[Strengthened scale selection]\label{lem:selection}
If $b$ has no positive power bound near zero, there are
$\gamma_i\downarrow0$, $s_i\to0$, actual bad parameters
$\lambda_i$, and associated $(U_i,z_i,r_i)$ with
\begin{align}
 \frac{b(s)}{s^{\gamma_i}}&\leq
 2\frac{b(s_i)}{s_i^{\gamma_i}}
 \quad(s_i\leq s\leq1/8),&
 \frac{b(s_i)}{s_i^{\gamma_i}}&>i,\label{eq:running-max}\\
 b(s_i)/4&<\lambda_i\leq b(s_i),&
 r_i&\leq s_iD(z_i).\label{eq:bad-choice}
 \intertext{For every fixed positive integer $k$, these satisfy}
 s_i/\lambda_i^k&\longrightarrow0.\label{eq:scale-separation}
\end{align}
\end{lemma}
\begin{proof}
Take $\gamma_i=1/i$. Failure of every power bound makes
$b(s)/s^{\gamma_i}$ unbounded near zero. Choose a truncated interval
$[a_i,1/8]$ whose finite supremum $M_i$ exceeds $2i$, and select
$s_i$ with quotient greater than $M_i/2$. This gives
\eqref{eq:running-max}; since $b\leq1$,
$s_i^{\gamma_i}<1/i$ and $s_i\to0$.
Choose an actual $\beta>b(s_i)/2$ and then an actual bad parameter
$\lambda_i>\beta/2$. This proves \eqref{eq:bad-choice} using only
approximate suprema. Finally
\[
 \frac{s_i}{\lambda_i^k}<(4/i)^k s_i^{1-k\gamma_i}\longrightarrow0.
\]
\end{proof}

Suppose now that such a bad sequence exists. If its $p_i$ have a
bounded subsequence, Lemma~\ref{lem:compact-p} gives a uniform Taylor
error $Cr^{1+\theta}$. The nonzero derivative $(Du,1)$ then makes
every $\lambda>Cr^\theta$ good. Thus
$\lambda_i\leq C's_i^\theta$, contradicting
$\lambda_i>i s_i^{\gamma_i}/4$. Hence we may assume $p_i\to\infty$
and $\epsilon_i\to0$.

For every fixed multiple of $r_i$ and every center within a fixed
multiple of $r_i$ of $z_i$, the parameter $\Lambda\lambda_i$,
$\Lambda=32$, is eventually good. Indeed for each fixed large $M$,
the boundary distances are comparable and
\[
 \beta\leq b(2Ms_i)\leq8\lambda_i(2M)^{\gamma_i}<32\lambda_i.
\]
Take a base fit $P_i$, rotate it to $c_ie_N$, and define
\begin{equation}\label{eq:anisotropic-scale}
 U_i(z_i+r_i(\sqrt{\lambda_i}x,t))
   =U_i(z_i)+r_ic_i[t+\lambda_i V_i(x,t)].
\end{equation}
The domains exhaust $\R^n\times\R$ by
\eqref{eq:scale-separation}. The following geometric conclusions of
\cite[Section 3]{OpenAI} apply with \eqref{eq:slope-bounds}:
the $V_i$ have a locally uniformly convergent subsequence with
nonaffine limit $V$, and $V$ obeys the all-center, all-radius fit
condition
\begin{equation}\label{eq:limit-fit}
 |V(x+h,t+\ell)-V(x,t)-p\cdot h-q\ell|\leq C_0R
 \quad(|\ell|\leq R,\ |h-p\ell|\leq R),
 \qquad C_0=8\Lambda.
\end{equation}
Here $p,q$ may depend on the center and $R$; writing
$q=a-|p|^2/2$ gives exactly the convention in
Theorem~\ref{thm:viscous}.

We spell out the reused compactness interfaces. Equal-parameter
concentric fits at comparable radii have direction and vector
differences $O(\sqrt\lambda)$ and norm difference $O(\lambda)$:
evaluate their affine errors on a common ball, then on a short common
axial segment. The displaced-center version holds in quarter margins.
Chaining from $r_i$ to $Rr_i$ costs a factor $1+\log R$. Choose a
fixed enclosing $R$ larger than the resulting multiple of
$1+\log R$ for each normalized compact box. In the base coordinates,
fit coefficients at its centers have
$P_x=O(\sqrt{\lambda_i})$ and $P_t-c_i=O(\lambda_i)$.
Thus the normalized coefficients
$p_i'=P_x/(c_i\sqrt{\lambda_i})$ and
$q_i'=(P_t/c_i-1)/\lambda_i$ are bounded. Their fitted domains
contain
\[
 |\ell|\leq S/2,\qquad
 |h-p_i'\ell/(1+\lambda_iq_i')|\leq S,
\]
with value error at most $2\Lambda S$. First choose $S$ small and
then increments small relative to the bounded coefficients; this
proves equicontinuity without derivative bounds on $V_i$. The
enclosing fit gives local boundedness. At each fixed center and
radius, take $S=4R$, extract the bounded fit coefficients, and use
strict domain margins to obtain \eqref{eq:limit-fit}. This auxiliary
subsequence may depend on that center and radius, since $V$ is
already fixed. Finally an affine limit $V=p\cdot x+qt$ would give
physical nonzero slopes $c_i(\sqrt{\lambda_i}p,1+\lambda_iq)$
with residual $o(\lambda_ir_i)$ on their $\lambda_i$ cylinders,
contradicting badness. Only the geometric part of the source's
extraction is used; its PDE identification is replaced next.

\section{The constrained equation and finite effective viscosity}

In the rotated frame put $\xi_i=(a_i,b_i)$,
$|a_i|^2+b_i^2=1$, $\beta_i=1/c_i$, $A_i=1-\beta_i^2$,
$\Pi_i=I-\xi_i\xi_i^T$ and
$S_\lambda=\operatorname{diag}(I,\sqrt\lambda)$.
At every test $\phi$ of $V_i$, the physical gradient is $c_iq_i$,
where
\[
 q_i=(\sqrt{\lambda_i}D_x\phi,1+\lambda_i\phi_t),
 \qquad q_i\cdot\xi_i=\beta_i.
\]
The last equality follows from the exact affine increment and holds
for both upper and lower tests. Substituting into \eqref{eq:lift}
and dividing by $\lambda_ic_i^3/r_i$ gives
\begin{align}
 D^2\phi[(D_x\phi,1+\lambda_i\phi_t)]
   +C_i(D\phi)\tr(S_{\lambda_i}\Pi_iS_{\lambda_i}D^2\phi)&=0,
 \label{eq:scaled-pde}\\
 C_i(q)=\nu_i+\epsilon_i(2q_t+|q_x|^2)
                   +\epsilon_i\lambda_iq_t^2,\qquad
 \nu_i&=\frac{\epsilon_iA_i}{\lambda_i}.\nonumber
\end{align}
The constraint was used in
\[
 |\Pi_iq_i|^2=A_i+\lambda_i(2\phi_t+|D_x\phi|^2)
                              +\lambda_i^2\phi_t^2.
\]
This identity is only asserted at actual tests. There,
$C_i=\epsilon_i|\Pi_iq_i|^2/\lambda_i\geq0$. Replacing $C_i$
globally by $\max\{C_i,0\}$ gives a continuous degenerate-elliptic
operator agreeing at every test; this legitimizes off-constraint
viscosity stability.

One bounded contact controls the possibly negative parameter $A_i$.
If $|V_i|\leq M$ on a fixed unit ball, then
$V_i-K|z|^2$, $K>2M+1$, attains its maximum in the interior.
Its quadratic upper test has uniformly bounded gradient. The
preceding nonnegative identity at that point yields
$A_i\geq-C_M\lambda_i$, and hence
\begin{equation}\label{eq:nu-lower}
 \nu_i\geq-C_M\epsilon_i\longrightarrow0\quad\hbox{from below}.
\end{equation}
No bounded test gradient is asserted at arbitrary points.

The exact increment also yields
\[
 V_i(z+sw_i)=V_i(z)+s\delta_i,\quad
 w_i=(a_i,b_i\sqrt{\lambda_i}),\quad
 \delta_i=\frac{\beta_i-b_i}{\sqrt{\lambda_i}}.
\]
Evaluating at $s=1/|w_i|$ in a larger fixed ball gives
$|\delta_i|\leq C|w_i|$. Pass to subsequences
$a_i\to a$, $b_i\to b$, $\delta_i\to\delta$ to obtain
\begin{equation}\label{eq:transport}
 V(x+sa,t)=V(x,t)+s\delta.
\end{equation}
If $a=0$, then $\delta=0$.

For a bounded subsequence of $\nu_i$,
\eqref{eq:nu-lower} gives $\nu_i\to\nu\geq0$. The extended
operators converge on compact jet sets, so their limit solves
\[
 D^2V[(D_xV,1)]+\nu\tr((I-aa^T)D^2_{xx}V)=0.
\]
If $a=0$, this is $G_\nu=0$. The fit condition
\eqref{eq:limit-fit}, the pinned Liouville theorem for $\nu=0$,
and Theorem~\ref{thm:viscous} for $\nu>0$ make $V$ affine.

If $a\neq0$, let $e=a/|a|$, $k=\delta/|a|$. The transport identity
gives $V(se+y,t)=ks+W(y,t)$, $y\perp e$. Tests $ks+\psi(y,t)$
are genuine tests for $V$, with Hessian annihilating $e$. Thus
\[
 D^2W[(D_yW,1)]+\nu\Delta_yW=0.
\]
The fits project with exactly the same constant: write an original
fit $p=de+p_\perp$ with ordinary time coefficient $q$ and choose
$h=d\ell e+h_\perp$. For $|\ell|\leq R$ and
$|h_\perp-p_\perp\ell|\leq R$ these increments lie in its box.
Subtracting the exact $kd\ell$ term gives
\[
 |W(y+h_\perp,t+\ell)-W(y,t)-p_\perp\cdot h_\perp
                 -(q+d^2-kd)\ell|\leq C_0R.
\]
The same rigidity applies in the $n-1\geq1$ transverse dimensions.
This again makes $V$ affine, contradicting its construction. All
finite-viscosity alternatives, including near-critical original
slopes, are excluded.

\section{A uniform small-solution estimate for the lift}

\begin{lemma}[Uniform small-solution interface]\label{lem:savin-family}
Fix $N$, $\lambda_0>0$, $\Lambda_0<\infty$, $\rho>0$ and $K<\infty$.
Consider autonomous continuous degenerate-elliptic operators $F(X,q)$
with $F(0,q)=0$. Suppose they are $C^2$ on the common jet ball
$|X|,|q|<\rho$, with common $C^2$ bound $K$, and their matrix
derivatives have eigenvalues in $[\lambda_0,\Lambda_0]$ there.
There are common $\sigma>0$, $\alpha>0$ and $C<\infty$ such that
any continuous viscosity solution of $F(D^2Z,DZ)=0$ on $B_1$ with
$\|Z\|_{L^\infty}\leq\sigma$ is $C^{2,\alpha}$ on $B_{1/2}$ with
norm at most $C$. The threshold can be reduced to obtain any fixed
sufficiently small bound on the first and second derivatives.
\end{lemma}
\begin{proof}
This is the common-data form of \cite[Theorem 1.3 and
Proposition 4.1, pp. 571--573]{Savin}. The proposition's contradiction
argument explicitly allows varying $F_k$ under a common second-derivative
bound, extracts $F_k,DF_k$, and obtains a uniformly elliptic
constant-coefficient limiting equation. Its quadratic approximation
and iteration constants depend only on the common local data.
The lower-order and space variables in that source are absent here.
Its H2 uses the matrix operator norm; the eigenvalue assumptions above
give its constants $\lambda_0$ and $N\Lambda_0$. Reducing the common
jet radius gives the last assertion. No gradient bound for the unknown
solution is assumed.
\end{proof}

\begin{lemma}[Coarse-cylinder estimate]\label{lem:coarse-cylinder}
Fix $n,L,b_0>0$. There are $d_*,k_*>0$ and $C<\infty$, depending
only on these data, with the following property. Suppose an actual
affine lift $U$ of a weak $p$-harmonic solution, $p\geq3$,
is defined on a neighborhood of an $\eta$-fitted cylinder
$Z_\eta(z,H;P)$, with slope $P=ce_N$ after rotation, and
\begin{equation}\label{eq:coarse-hyp}
 \begin{gathered}
 c\in[1/2,L+1],\quad A=1-c^{-2}>0,\quad
 \epsilon=(p-2)^{-1}\leq1,\quad e=\epsilon A\leq1,\\
 \xi=(a,b),\quad |a|^2+b^2=1,\quad b\geq b_0,\qquad 0<\eta\leq d_*e.
 \end{gathered}
\end{equation}
Then $P_{\mathrm{new}}=DU(z)$ exists and is nonzero, and, whenever
$0<\lambda\leq e$, $0<r\leq k_*\ell$ with
$\ell=(H/4)\sqrt{\eta/e}$, one has
\begin{equation}\label{eq:cylinder-estimate}
 \sup_{h:\,z+h\in Z_\lambda(z,r;P_{\mathrm{new}})}
 |U(z+h)-U(z)-P_{\mathrm{new}}\cdot h|\leq Ce r^2/H.
\end{equation}
\end{lemma}
\begin{proof}
Put $T_e=\operatorname{diag}(\sqrt e I,1)$ and set
\[
 U(z+\ell T_e(x,t))=U(z)+c\ell[t+eW(x,t)].
\]
For $d_*\leq1$ this ellipsoid lies strictly in the fitted cylinder:
its transverse radius is $H\sqrt\eta/4$ and its axial radius is at
most $H/4$. The value fit gives
\begin{equation}\label{eq:W-small}
 \|W\|_{L^\infty(B_1)}\leq(4/c)\sqrt{\eta/e}.
\end{equation}
At constrained tests its equation is
\[
 D^2W[(D_xW,1+eW_t)]+B(DW)\tr(S_e\Pi S_eD^2W)=0,
\]
where $S_e=\operatorname{diag}(I,\sqrt e)$ and
\[
 B(q)=1+\epsilon(2q_t+|q_x|^2)+\epsilon^2Aq_t^2.
\]
At those tests $B=|\Pi(\sqrt e q_x,1+e q_t)|^2/A\geq0$.
Again replace $B$ globally by its nonnegative part. To obtain a
proportional, rather than absolute, derivative bound include an
amplitude parameter $d\in[0,1]$:
\begin{align*}
 F_d(X,q)&=X[(d q_x,1+edq_t)]
           +\max\{B_d(q),0\}\tr(S_e\Pi S_eX),\\
 B_d(q)&=1+\epsilon(2dq_t+d^2|q_x|^2)
                          +\epsilon^2Ad^2q_t^2.
\end{align*}
Allowing $\epsilon,A,e,d$ independently in $[0,1]$ only enlarges
the family. All members are continuous and degenerate elliptic on
all jets. On $|q|\leq1/8$ one has $B_d\geq3/4$, so the max switch
is absent. They vanish at $X=0$ and have common local $C^2$ bounds.
At $q=0$ their matrix form is
\[
 Q(h,s)=|h|^2+(1+e)s^2-(a\cdot h+b\sqrt e s)^2,
 \qquad \frac{b_0^2}{3}(|h|^2+s^2)\leq Q(h,s)
                    \leq2(|h|^2+s^2).
\]
Here is a dimension-free certificate for this actual vector form.
Set $t=\sqrt e$ and
$g=|a|^2|h|^2-\langle a,h\rangle^2\geq0$, by Cauchy--Schwarz.
The unit relation $|a|^2+b^2=1$ gives
\[
 Q=|bh-tsa|^2+s^2+g
\]
and
\[
 \begin{split}
 3Q-b^2(|h|^2+s^2)
 ={}&|bh-2tsa|^2+|bh-tsa|^2+3g\\
    &+(3-b^2-2t^2|a|^2)s^2\geq0.
 \end{split}
\]
Indeed $3-b^2-2t^2|a|^2\geq1+b^2$ for $t^2\leq1$.
The upper bound follows by dropping the negative square in $Q$.
This works on any real inner product space and includes $a=0$.
The supporting Lean theorem \texttt{physical\_uniform\_bounds}
certifies precisely these bounds with $t=\sqrt e$, the unit relation,
$0\leq e\leq1$ and $b_0^2\leq b^2$.

We quantify the common neighborhood as well. Write
\[
 v(q)=(dq_x,1+edq_t),\qquad
 M_d(q)=v(q)v(q)^T+B_d(q)S_e\Pi S_e.
\]
For $|q|\leq r\leq1/8$, one has $|v(q)-e_N|\leq r$,
$|B_d(q)-1|\leq2r+2r^2$, and $\|S_e\Pi S_e\|\leq1$. Hence
\[
 \|M_d(q)-M_d(0)\|\leq4r+3r^2<5r.
\]
Thus the explicit common radius
\[
 \rho=\min\{1/8,b_0^2/30\}
\]
gives $B_d\geq3/4$ and eigenvalues in $[b_0^2/6,3]$ throughout
$|q|\leq\rho$, including all parameter endpoints. For every positive
semidefinite matrix $Y$, the operator-norm convention gives
\[
 \frac{b_0^2}{6}\|Y\|\leq\tr(M_d(q)Y)
                  \leq3\tr Y\leq3N\|Y\|.
\]
Savin's ellipticity constants can therefore be taken as
$\lambda=b_0^2/6$ and $\Lambda=3N$. On $|X|,|q|\leq\rho$ the
max switch is absent, so the coefficients are polynomial in $q$;
$D_qM_d$ and $D_q^2M_d$ are uniformly bounded. Consequently
$D^2_{Xq}F_d=D_qM_d$ and
$D^2_{qq}F_d=\tr(D_q^2M_dX)$ have common bounds depending only on
$N,b_0$. Together with $F_d(0,q)=0$, these verify every hypothesis
of Lemma~\ref{lem:savin-family} with one common $K$. The square
root in $S_e$ causes no loss: parameter derivatives are not required.

Let $a_0=\|W\|\leq\sigma$, with the common threshold from
Lemma~\ref{lem:savin-family}. If $a_0=0$ the estimate is immediate.
Otherwise take $d=a_0/\sigma$ and $Z=W/d$. Dividing the equation by
$d$ gives exactly $F_d(D^2Z,DZ)=0$, and $\|Z\|=\sigma$. Thus
\[
 \|DW\|_{B_{1/2}}+\|D^2W\|_{B_{1/2}}\leq C a_0.
\]
Choose $d_*$ so that \eqref{eq:W-small} is below the threshold.
Scaling back yields
\[
 \|T_eD^2UT_e\|\leq\frac{ce}{\ell}C\frac4c\sqrt{\eta/e}
                  \leq C'e/H
\]
on the half ellipsoid. By decreasing the threshold its derivative
conclusion also gives $c(1+eW_t)\geq c/2$. Hence
\[
 P_{\mathrm{new}}=c(\sqrt e D_xW(0),1+eW_t(0))\neq0,
 \qquad |T_e^{-1}N_{P_{\mathrm{new}}}|\leq C.
\]
For $h$ in its $\lambda$ cylinder, $\lambda\leq e$ implies
$|T_e^{-1}h|\leq Cr+\sqrt{\lambda/e}r\leq C'r$.
Choose $k_*$ so the whole segment lies in the half ellipsoid.
Taylor's formula proves \eqref{eq:cylinder-estimate}.
\end{proof}

\section{Infinite effective viscosity and the quadratic comparison}

It remains to exclude $\nu_i\to+\infty$. Then $A_i>0$ eventually.
Define
\[
 R_i=\nu_i^2,\quad H_i=R_ir_i,\quad
 \kappa_i=64(4R_i)^{\gamma_i},\quad \eta_i=\kappa_i\lambda_i.
\]
Since $A_i,\epsilon_i\leq1$, $R_i\leq\lambda_i^{-2}$, and
\eqref{eq:scale-separation} gives $R_is_i\to0$. Centers within
$2H_i$ of $z_i$ retain at least half its boundary distance. At all
admissible radii up to $H_i$, \eqref{eq:running-max} gives
\[
 \beta\leq b(4R_is_i)\leq8\lambda_i(4R_i)^{\gamma_i}<\eta_i.
\]
Thus $\eta_i$ is good throughout these comparisons. Also
$\eta_i\leq C\lambda_i^{1-2\gamma_i}\to0$ and
$\kappa_i/\nu_i\to0$. A fixed multiple of $\lambda_i$ is not
asserted good on these growing boxes.

Let $Q_{i,0}$ be an $\eta_i$-fit at the base radius $r_i$, and
let $P_i$ be the original $\Lambda\lambda_i$-fit. Their common
inscribed ball gives
$|P_i-Q_{i,0}|\leq C\kappa_i\sqrt{\lambda_i}$, hence the same
order for their axis angle. An equal-parameter norm comparison
cannot be applied to this pair. Instead take a point on the base
axis at distance $t_ir_i$ where $t_i=c\kappa_i^{-1/2}$ and $c>0$
is sufficiently small and fixed. Its transverse displacement from
the coarse axis is at most
$Ct_i\kappa_i\sqrt{\lambda_i}r_i\leq\sqrt{\eta_i}r_i$;
the axial margins also hold. Evaluating both fits there gives
\[
 |(P_i-Q_{i,0})\cdot N_{P_i}|
       \leq C\kappa_i^{3/2}\lambda_i.
\]
Converting axial projection to norm costs the square of the angle,
at most $C\kappa_i^2\lambda_i$. Therefore
\begin{equation}\label{eq:unequal-norm}
 \bigl||P_i|-|Q_{i,0}|\bigr|\leq C\kappa_i^2\lambda_i.
\end{equation}
Now chain $\eta_i$-fits through dyadic radii from $r_i$ to $H_i$.
These comparisons have equal parameters and common slope bounds,
so the source's quadratic norm comparison is applicable. A final
coarse fit $Q_i$ satisfies
\begin{equation}\label{eq:coarse-norm}
 \bigl||Q_i|-c_i\bigr|\leq
 C\{\kappa_i^2\lambda_i+\eta_i(1+\log\nu_i)\}.
\end{equation}
Put $c_i'=|Q_i|$, $A_i'=1-(c_i')^{-2}$. The map
$c\mapsto1-c^{-2}$ is uniformly Lipschitz on $[1/2,L+1]$, so
\[
 \frac{|A_i'-A_i|}{A_i}
 \leq C\epsilon_i
   \frac{\kappa_i^2+\kappa_i(1+\log\nu_i)}{\nu_i}
 \longrightarrow0.
\]
Indeed $\kappa_i=64\,4^{\gamma_i}\nu_i^{2\gamma_i}$,
$\gamma_i\to0$ and $\nu_i\to\infty$; eventually
$4\gamma_i<1/2$. Thus $A_i'>0$ and
$e_i'=\epsilon_i A_i'\sim\nu_i\lambda_i$.
This proves the needed comparison before substituting the coarse
parameter into the small-solution lemma.

The coarse fit has
$\eta_i/e_i'\sim\kappa_i/\nu_i\to0$. Its added-coordinate
component obeys
$b_i'\geq(1-\sqrt{\eta_i})/(L+1)$ by
\eqref{eq:slope-bounds}. Its ellipsoid scale is
\[
 \frac{\ell_i}{r_i}
 =\frac{R_i}{4}\sqrt{\eta_i/e_i'}
 \sim\frac14\nu_i^{3/2}\sqrt{\kappa_i}\longrightarrow\infty.
\]
Also $\lambda_i\leq e_i'$ eventually. Estimate
\eqref{eq:cylinder-estimate} at $r_i$ therefore has error, divided
by $\lambda_ir_i$, bounded by
\[
 \frac{Ce_i'r_i}{\lambda_iH_i}\sim\frac C{\nu_i}\longrightarrow0.
\]
Its nonzero slope is the actual $DU_i(z_i)$. It makes
$\lambda_i$ good, contradicting its selection. Every domain used
lies in the original solution region by
\eqref{eq:scale-separation}. We have excluded all branches of the
bad sequence and proved the required power bound.


\section{Gradient extraction and a consequence}

The preceding argument yields one power bound $b(s)\leq Bs^\gamma$
for the full family $p\geq3$. Decrease $\gamma$ so $0<\gamma\leq1$.
At centers in $B_{1/2}\subset\R^{n+1}$ choose common $A_0,r_0>0$
so $\lambda=A_0r^\gamma$ is good for $0<r\leq r_0$.
On dyadic radii $r_j=2^{-j}r_0$, the inscribed balls have radii
$\rho_j=\sqrt{A_0}r_j^{1+\gamma/2}$ and error
$A_0r_j^{1+\gamma}$. Comparing consecutive affine functions on a
common ball gives
\[
 |P_j(z)-P_{j+1}(z)|\leq Cr_j^{\gamma/2}.
\]
The slopes converge uniformly. For
$\rho_{j+1}<|h|\leq\rho_j$, the residual divided by $|h|$ is at most
$Cr_j^{\gamma/2}\leq C'|h|^\alpha$, where
$\alpha=\gamma/(2+\gamma)\in(0,1/3]$.
Thus $U$ is differentiable with uniform Taylor error
$C|h|^{1+\alpha}$. Comparing these Taylor formulas at two centers,
at their connecting point and at its translate by $|z-z'|e$ for each
unit vector $e$, gives the $\alpha$-Hölder bound on $DU$. Larger
separations use the Lipschitz bound. Restrict to $b=0$ and restore
the original oscillation to prove Theorem~\ref{thm:main}.

If $u_p$ converges locally uniformly in value to $v$ with bounded
local oscillations, the same Taylor estimate passes to $v$: at each
fixed center, take a convergent subsequence of its bounded slopes.
The limiting Taylor inequality identifies a unique derivative with
that error. The two-center comparison gives the same exponent for
$v$; it need not agree with the independently supplied
infinity-harmonic exponent. On nested balls, finite differences give
\[
 \|Du_p-Dv\|\leq C(h^\alpha+\delta/h),
 \qquad\delta=\|u_p-v\|_{L^\infty}.
\]
For small $\delta$, choose $h=\delta^{1/(1+\alpha)}$ to obtain
$\|Du_p-Dv\|\leq C\delta^{\alpha/(1+\alpha)}$. A rate in $p$ requires
independent value-convergence assumptions.

\section{Verification and scope}

The proof relies on the pinned mathematical inputs listed in Section 1.
It gives no explicit numerical exponent, boundary estimate, endpoint
exponent, or result for the inhomogeneous equation beyond the stated
large-$p$ corollary. The literature comparison is bounded; no worldwide
priority claim is made.

\begin{samepage}
The accompanying Lean 4.34.1 project is pinned to mathlib commit
\begin{center}\texttt{d13f23b723b8a846827a245b89c10fc7d3f11612}.\end{center}
\end{samepage}
It proves ten supporting theorems: seven vector inner-product and
ellipticity theorems and the original three scalar real-algebra
lemmas. The vector certificate applies in any real inner product
space. Its \texttt{physical\_uniform\_bounds} theorem directly bounds
\[
 |h|^2+(1+e)s^2-(\langle a,h\rangle+b\sqrt e\,s)^2
\]
under $|a|^2+b^2=1$, $0\leq e\leq1$ and $b_0^2\leq b^2$,
without a coordinate split. The compiled theorems have only the
standard Lean dependencies \texttt{propext},
\texttt{Classical.choice} and \texttt{Quot.sound}. They do not
formalize the local jet-neighborhood argument, Savin's theorem,
the PDE, geometric compactness, source theorems, or the main
conclusion. Exact arithmetic scripts check
selected identities and reject specified mutations; they do not
establish analytic or universal claims. The reproducibility receipt
separates these coverage levels.

\appendix
% Publication adaptation of I236, frozen SHA256 a30077ca36f40e6613e340d3f8ae0ac692000f7b39d6eb4113897f40b7c68931.
% New presentation requires its own review; this file is not the reviewed original.
\section{Finite-viscosity fitted-class rigidity}\label{app:viscous}

This appendix supplies the positive-viscosity input used in the main
argument. Its proof of record is I236, \emph{analytic-attempt},
Sections 1--6, previously reviewed in I239. The full argument is
carried here rather than replaced by that review verdict.
It includes the contact-propagation proof; it does not assume
regularity of arbitrary entire fitted solutions. Its geometric inputs
are the fit compactness and barycentric-envelope lemmas of
\cite[Sections 4--5]{OpenAI}, with the equation changes checked below.

Fix an integer $m\geq1$ and $0<\tau<\infty$. For $P=(p,q)\in
\mathbb R^m\times\mathbb R$ and a symmetric matrix $X$, define
\[
 A_\tau(p)=\begin{pmatrix}\tau I+pp^T&p\\p^T&1\end{pmatrix},
 \qquad G_\tau(P,X)=\operatorname{tr}(A_\tau(p)X).
\]
The time-gradient component $q$ does not enter $G_\tau$. We use the
positive-Hessian convention: upper tests for subsolutions give
$G_\tau\geq0$, and lower tests for supersolutions give $G_\tau\leq0$.

\begin{theorem}[Finite-viscosity rigidity]\label{thm:viscous}
Let $V\in C(\mathbb R^m\times\mathbb R)$ solve $G_\tau=0$ in the
viscosity sense. Suppose a finite $C_0\geq0$ has the following property:
for every $(x,t)$ and every $R>0$, some $p\in\mathbb R^m$ and
$a\in\mathbb R$ satisfy
\begin{equation}\label{eq:app-fit}
 \left|V(x+h,t+\ell)-V(x,t)-p\cdot h
       -\left(a-\frac{|p|^2}{2}\right)\ell\right|\leq C_0R
 \quad (|\ell|\leq R,\ |h-p\ell|\leq R).
\end{equation}
Then $V$ is affine. No gradient bound or energy assumption is required.
No estimate uniform as $\tau$ approaches zero or infinity is asserted.
\end{theorem}

\subsection{Time regularity of spatially semiconcave subsolutions}

\begin{lemma}\label{lem:time-lipschitz}
A continuous $G_\tau$-subsolution with locally uniform spatial
semiconcavity is locally Lipschitz in all variables.
\end{lemma}
\begin{proof}
Work on a product cylinder compactly contained in the domain and a
smaller cylinder of base points $(x_0,t_0)$. Choose $r>0$ so that every
$\overline B_r(x_0)$ stays inside the larger spatial ball. Enlarge the
semiconcavity constant to $K\geq0$. Continuity supplies a common bound
$M$ and a time modulus $\omega$ on the larger compact cylinder.

At every base point, a supporting hyperplane to
$u(\cdot,t_0)-K|\cdot|^2/2$ gives a vector $p$ with
\[
 u(y,t_0)\leq u(x_0,t_0)+p\cdot(y-x_0)+\frac K2|y-x_0|^2
 \quad (|y-x_0|\leq r).
\]
Evaluating at opposite fixed increments bounds $|p|$ by a common $L$.
The same argument at neighboring points bounds spatial Lipschitz
constants uniformly in time. Set $K'=K+1$ and hold this $p$ fixed in
\[
 h(t)=\max_{|y-x_0|\leq r}
 \left\{u(y,t)-p\cdot(y-x_0)-\frac{K'}2|y-x_0|^2\right\}.
\]
At $t_0$ the unique maximizer is $x_0$, and boundary competitors have
a deficit at least $r^2/2$. Choose a common $s>0$ with
$\omega(s)<r^2/8$ and with the relevant time intervals inside the
larger cylinder. For $|t-t_0|<s$, boundary competitors are at most
$u(x_0,t_0)-3r^2/8$, whereas the central competitor is at least
$u(x_0,t_0)-r^2/8$. Thus every maximizing point is interior. Also
$h$ is continuous and $|h|\leq H:=M+Lr+K'r^2/2$.

If a scalar upper test $\psi$ touches $h$ at $t$ and $y$ is a maximizer,
then
\[
 \phi(Y,T)=\psi(T)+p\cdot(Y-x_0)+\frac{K'}2|Y-x_0|^2
\]
is an upper test of $u$ at $(y,t)$, after adjusting the contact constant.
Its mixed Hessian is zero. The subsolution inequality gives
\[
 0\leq\psi''(t)+K'\bigl(|p+K'(y-x_0)|^2+\tau m\bigr),
 \qquad \psi''(t)\geq-C,
 \quad C=K'\bigl((L+K'r)^2+\tau m\bigr).
\]
Consequently $h(t)+C(t-t_0)^2/2$ is convex. The one-dimensional
viscosity characterization used here follows directly from the chord
criterion: failure of convexity yields an interior upper contact by a
strictly concave quadratic added to an affine function.
The slopes of this bounded convex function are bounded on the middle
half interval by its endpoint secants. For example, $8H/s+4Cs$ bounds
the Lipschitz constant of $h$ there. Since
$u(x_0,t)\leq h(t)$ with equality at $t_0$, this bounds
$u(x_0,t)-u(x_0,t_0)$ from above. Repeating at base time $t$, with the
same constants, gives the reverse bound. Combining with the spatial
bound proves the lemma.
\end{proof}

\subsection{Preserving contact while making the supersolution Lipschitz}

Suppose $u\leq v$, $u$ is locally Lipschitz, $v$ is a continuous
supersolution, and $u(z_0)=v(z_0)$. Work inside a ball with a compact
larger margin. If $L_0$ bounds the Lipschitz constant of $u$ there,
choose $L>L_0$ and a small $r>0$. Define locally
\[
 v_L(z)=\min_{|h|\leq r}\{v(z+h)+L|h|\}.
\]
The minimum is attained, and $v_L$ is continuous. Every competitor
obeys
\begin{equation}\label{eq:cone-lower}
 v(z+h)+L|h|\geq u(z)+(L-L_0)|h|.
\end{equation}
Thus $u\leq v_L\leq v$, and the unique minimizer at $z_0$ is zero.
For any minimizing $h$,
\[
 (L-L_0)|h|\leq v(z)-u(z).
\]
Continuity therefore places all minimizers in $B_{r/2}$ on a small
neighborhood of $z_0$.

A lower test $\phi$ of $v_L$ at $z$, pulled back by a minimizing
fixed translation $h$, is a lower test of $v$ at $z+h$, with unchanged
gradient and Hessian. The autonomous equation implies that $v_L$
remains a supersolution. If $h$ minimizes at $z$, then
$h+z-z'$ is admissible at a sufficiently close $z'$, so
\[
 v_L(z')\leq v(z+h)+L|h+z-z'|
           \leq v_L(z)+L|z-z'|.
\]
Reverse the roles of the two points to obtain local $L$-Lipschitz
continuity. Finally, if $v_L(z)=u(z)$, equation
\eqref{eq:cone-lower} forces its minimizing shift to be zero, whence
$v(z)=u(z)$. This last implication recovers contact for the original
supersolution.

\subsection{The exponential flux and weak inequalities}

On the full gradient space put
\[
 H(p,q)=q+\frac{|p|^2}{2},\qquad
 \mathcal A(p,q)=e^{H(p,q)/\tau}(p,1).
\]
Direct differentiation gives
\begin{equation}\label{eq:exp-flux}
 D\mathcal A(P)=\frac{e^{H(P)/\tau}}{\tau}A_\tau(p),
 \qquad
 \operatorname{div}\mathcal A(Df)
       =\frac{e^{H(Df)/\tau}}{\tau}G_\tau(Df,D^2f)
\end{equation}
for smooth $f$. This exponential structure is classical; see
\cite{SanchezMorgado}. It is not a new operator introduced here.
The derivative is symmetric positive definite. On $|P|\leq M$, one
may take
\[
 \begin{aligned}
 \lambda_0&=\bigl[\max\{2,(1+2M^2)/\tau\}\bigr]^{-1},
 &\Lambda_0&=M^2+1+\tau,\\
 \lambda_M&=e^{-M/\tau}\lambda_0/\tau,
 &\Lambda_M&=e^{(M+M^2/2)/\tau}\Lambda_0/\tau
 \end{aligned}
\]
in $\lambda_MI\leq D\mathcal A(P)\leq\Lambda_MI$.
Indeed,
$A_\tau(p)[(\zeta,s)]=(p\cdot\zeta+s)^2+\tau|\zeta|^2$;
the lower bound follows from
$s^2\leq2(s+p\cdot\zeta)^2+2M^2|\zeta|^2$.

\begin{lemma}\label{lem:flux-weak}
For a locally Lipschitz viscosity subsolution $f$,
$\operatorname{div}\mathcal A(Df)\geq0$ distributionally.
For a locally Lipschitz supersolution the inequality is reversed.
\end{lemma}
\begin{proof}
Fix a compact test support in a region where $f$ is $M$-Lipschitz,
with an additional compact margin. For small $\epsilon>0$ take the
full-variable sup convolution
\[
 f^\epsilon(z)=\sup_y\left(f(y)-\frac{|z-y|^2}{2\epsilon}\right).
\]
Optimizers satisfy $|y-z|\leq2M\epsilon$. On the smaller region
$f^\epsilon$ is $M$-Lipschitz, converges uniformly to $f$, and has
$D^2f^\epsilon\geq-I/\epsilon$ distributionally. Pulling an upper
test to an attaining fixed translate shows that $f^\epsilon$ is an
exact subsolution. This uses autonomy, not convexity in the gradient.
At its almost every Alexandrov twice-differentiability point,
\[
 \operatorname{tr}\bigl(D\mathcal A(Df^\epsilon)D^2f^\epsilon\bigr)
 \geq0.
\]
To see this, add a positive arbitrarily small quadratic to its
second-order Taylor polynomial and use the viscosity inequality.

Mollify $f^\epsilon$ on a smaller region, with mollifications
$f_j^\epsilon$. Then $|Df_j^\epsilon|\leq M$ and
$D^2f_j^\epsilon\geq-I/\epsilon$, giving the common lower bound
\[
 \operatorname{div}\mathcal A(Df_j^\epsilon)
 \geq-(m+1)\Lambda_M/\epsilon.
\]
Both gradients and Hessians converge almost everywhere to those of
$f^\epsilon$. For the Hessian statement, the distributional Hessian
of a semiconvex function is a matrix measure; its absolutely continuous
part is the Alexandrov Hessian, and mollifications of the singular part
converge to zero almost everywhere. This is the mollification fact in
\cite[Remark 3.3]{JulinJuutinen}, with the sign reversed from its
semiconcave formulation. For every nonnegative smooth compactly
supported $\phi$, integration by parts and one-sided Fatou give
\[
 \begin{split}
 -\int\mathcal A(Df^\epsilon)\cdot D\phi
 &=\lim_j\int\phi\,\operatorname{div}\mathcal A(Df_j^\epsilon)\\
 &\geq\int\phi\,
 \operatorname{tr}\bigl(D\mathcal A(Df^\epsilon)D^2f^\epsilon\bigr)
 \geq0.
 \end{split}
\]
In particular the displayed nonnegative density is locally integrable.
The favorable singular Hessian contribution has not been deleted.

It remains to pass to the nonlinear flux as $\epsilon\downarrow0$.
At a differentiability point $z$ of $f$, write an optimizer as
$y=z+\epsilon b_\epsilon$, where $|b_\epsilon|\leq2M$.
Taylor expansion, uniform for bounded $b_\epsilon$, and comparison
with $y=z+\epsilon Df(z)$ give
\[
 \tfrac12|b_\epsilon-Df(z)|^2=o(1).
\]
At a differentiability point of $f^\epsilon$, its gradient is
$b_\epsilon$. Along a fixed sequence $\epsilon_j\downarrow0$,
discard the null sets for $f$ and all its convolutions. Gradients then
converge almost everywhere and are uniformly bounded, so the smooth
flux converges in $L^1$ on the test support. This proves the first
assertion. For supersolutions use inf convolution, whose Hessian has
the upper bound $I/\epsilon$. Apply Fatou to the negative divergence
and compare its minimizing shift with $-\epsilon Df(z)$. This gives
the reversed inequality, without claiming that $-f$ satisfies the
same equation.
\end{proof}

\subsection{Contact propagation}

\begin{theorem}\label{thm:contact}
Let $\Omega\subset\mathbb R^{m+1}$ be connected and open. Suppose
$u\leq v$ are continuous, $u$ is a $G_\tau$-subsolution, $v$ is a
$G_\tau$-supersolution, and $u$ is locally uniformly semiconcave in
space. If they agree at one interior point, they agree everywhere.
\end{theorem}
\begin{proof}
First let both functions be locally Lipschitz and put $d=v-u\geq0$.
On a region where their gradients have norm at most $M$, define
\[
 B(z)=\int_0^1D\mathcal A\bigl(Du(z)+s(Dv(z)-Du(z))\bigr)\,ds.
\]
This is measurable, symmetric and uniformly elliptic between
$\lambda_MI$ and $\Lambda_MI$. Lemma~\ref{lem:flux-weak} and the
fundamental theorem of calculus imply
\[
 \mathcal A(Dv)-\mathcal A(Du)=B Dd,\qquad
 \int B Dd\cdot D\phi\geq0\quad(\phi\geq0).
\]
At a zero $z_0$ of $d$, choose any sufficiently small ball $B_r(z_0)$
compactly inside this region. Lax--Milgram gives $w\in H^1(B_r)$ with
$w-d\in H^1_0(B_r)$ and
$\int BDw\cdot D\phi=0$ for $\phi\in H^1_0(B_r)$.
The negative-part test gives $w\geq0$, because the boundary trace is
nonnegative. Testing the difference with $(w-d)^+$ gives
\[
 \int BD(w-d)^+\cdot D(w-d)^+\leq0,
\]
so $0\leq w\leq d$. Classical interior continuity and Harnack for
linear divergence-form weak solutions with bounded measurable
symmetric uniformly elliptic coefficients apply
\cite[Sections 2--3]{Crescimbeni}. For $m=1$ the total dimension is
two. If using a version stated only in dimensions above two,
extend $w$ constantly in a third coordinate and extend the coefficient
matrix to $\operatorname{diag}(B,1)$ on a product neighborhood.
Fubini verifies the weak equation, with common positive ellipticity;
the three-dimensional Harnack inequality then restricts to the
original plane. Thus $w(z_0)=0$ and Harnack,
followed by a chain of balls, gives $w=0$ throughout $B_r$.
Since $w-d\in H^1_0(B_r)$, $d$ has zero trace there. Its Lipschitz
continuity on a neighborhood of the closure identifies that trace
with its pointwise boundary restriction. Hence $d=0$ on the entire
sphere of radius $r$. The radius was arbitrary below a fixed margin,
so $d$ vanishes on a neighborhood of $z_0$.

For the stated pair, Lemma~\ref{lem:time-lipschitz} first makes $u$
locally Lipschitz. At every contact point apply the preceding cone
construction to $v$. The Lipschitz-pair argument makes $v_L=u$
nearby, and \eqref{eq:cone-lower} recovers $v=u$ there. The contact
set is consequently open. It is closed by continuity, nonempty by
hypothesis, and all of $\Omega$ by connectedness.
\end{proof}

\subsection{Geometric normalization and the affine-slice alternative}

We now prove Theorem~\ref{thm:viscous}. If $C_0=0$, fitted open boxes
are exactly affine and overlap to give one affine function. Suppose
$C_0>0$; it may be enlarged to at least one for estimates.
For a fitted frame $(x_0,t_0,R,p_0,a_0)$ the transformation
\begin{equation}\label{eq:app-frame}
 U(y,s)=\frac{V(x_0+R(y+p_0s),t_0+Rs)-V(x_0,t_0)}R
        -p_0\cdot y-\left(a_0+\frac{|p_0|^2}{2}\right)s
\end{equation}
preserves the fit constant, changing fit parameters to
$(p-p_0,a-a_0)$. Differentiation on test functions gives
$G_\tau[U]=R G_\tau[V]$, with $\tau$ unchanged. The equation-free
fit comparisons and compactness in \cite[Section 4]{OpenAI} therefore
apply: normalized functions have centered growth at most
$CK(1+\log K)^2$ on boxes of radius $K$, and form a locally compact
family. Locally uniform limits of these normalized frames retain
both $G_\tau=0$ and the all-center, all-radius fit
\eqref{eq:app-fit} with the same $C_0$. The equation passes by
viscosity stability on finite jets. For any fixed center and radius,
fit compactness gives a bounded sequence of fit coefficients;
extract a convergent subsequence and use strict fitted-domain margins
to pass the fit inequality, then continuity to include the boundary.
The auxiliary subsequence may depend on that center and radius,
since the limiting function is already fixed. This proves the
required fit for every center and radius without choosing one
subsequence for uncountably many fits.
The growth estimate supplies uniform spatial subquadratic tails on
compact time intervals.

If every horizontal slice is affine, write
$V(x,t)=P(t)\cdot x+c(t)$, with $P,c$ continuous. A fit at $(x,t)$
of radius $R$ gives $|p-P(t)|\leq C_0$ by taking $\ell=0$. For fixed
$\ell>0$ and $R\geq\ell$, evaluate the same fit at
$(h,\ell)=(p\ell,\ell)$ and $(-p\ell,-\ell)$ and add to obtain
\[
 \begin{split}
 \big|&[P(t+\ell)+P(t-\ell)-2P(t)]\cdot x\\
 &+\ell p\cdot[P(t+\ell)-P(t-\ell)]
      +c(t+\ell)+c(t-\ell)-2c(t)\big|\leq2C_0R.
 \end{split}
\]
All terms except the coefficient of $x$ are bounded uniformly as $x$
varies. Thus that coefficient vanishes, and continuity makes
$P(t)=bt+d$. The centered growth bound, applied at $x=\pm Ke$ and
$|t|\leq K$, yields $|P(t)|\leq C(1+\log K)^2$ and forces $b=0$.
Tests $P\cdot x+\psi(t)$ then reduce the equation for $c$ to
$\psi''=0$, so $c$ and $V$ are affine.

Otherwise a horizontal slice has a nonzero midpoint defect. The
transformation $V(x,t)\mapsto-V(x,-t)$ preserves the equation and
the fits, changes $(p,a)$ to $(-p,a)$, and reverses the defect. Choose
its sign positive.

\subsection{Extremal envelope and the changed equation}

Let $\kappa>0$ be the supremum, over all times and finite
positive-weight probability configurations, of
\[
 \frac{\sum_j\theta_jV(y_j,t)-V(\sum_j\theta_jy_j,t)}r,
 \qquad
 r^2=\sum_j\theta_j\left|y_j-\sum_k\theta_ky_k\right|^2>0.
\]
It is finite: normalize a radius-$r$ fit at the barycenter using
\eqref{eq:app-frame}. The centered growth bound is at most
$C(1+|y|^2)$; integration against a mean-zero, unit-second-moment
configuration bounds its normalized defect. Both the quotient and
$\tau$ are invariant under the fitted-frame transformation.

Take configurations tending to $\kappa$, normalized to mean zero and
second moment one, and locally uniformly convergent normalized
solutions $V_j\to\widehat V$. The corresponding measures have a weak
limit $\mu$ with mean zero and second moment at most one. First-moment
uniform integrability preserves the mean; uniform subquadratic tails
give uniform integrability for the varying value integrals. Hence
\[
 \int\widehat V(y,0)\,d\mu(y)-\widehat V(0,0)=\kappa.
\]
For $s>0$ use the barycentric envelope
\[
 \mathcal C_sV(x,t)=
 \sup_{\substack{\theta_j>0,\ \sum\theta_j=1\\
                   \sum\theta_jy_j=x}}
 \left\{\sum_j\theta_jV(y_j,t)
              -\frac1{2s}\sum_j\theta_j|y_j-x|^2\right\}.
\]
The Jensen bound, also valid in the normalized closure, gives
\begin{equation}\label{eq:envelope-sandwich}
 0\leq\mathcal C_sV-V\leq
 \sup_{r\geq0}\left(\kappa r-\frac{r^2}{2s}\right)
 =\frac{\kappa^2s}{2}.
\end{equation}
The measure $\mu$ is admissible for the envelope by its geometric
measure representation \cite[Section 5]{OpenAI}. At $s=1/\kappa$
its defect gives the reverse bound $\kappa/2$. Equality also forces
the second moment of $\mu$ to be one. Finite attainment, continuity
and the concave-envelope identity are the geometric conclusions of
that same lemma. A radius-$1/\kappa$ frame and time translation put
the equality at parameter one and time one. Rename the resulting
solution $V$, and set on $t>0$
\[
 W(x,t)=\mathcal C_tV(x,t)-\frac{\kappa^2t}{2}.
\]
Then
\begin{equation}\label{eq:extremal-contact}
 V-\frac{\kappa^2t}{2}\leq W\leq V,\qquad
 W(0,1)=V(0,1)=0,\qquad
 \mathcal C_1V(0,1)-V(0,1)=\frac{\kappa^2}{2}>0.
\end{equation}
The geometric identity
\[
 \mathcal C_tV(x,t)=\frac{|x|^2}{2t}
  +\operatorname{Conc}\left[V(\cdot,t)-\frac{|\cdot|^2}{2t}\right](x)
\]
gives $D^2_{xx}W\leq I/t$. Here $\operatorname{Conc}$ denotes the
least concave majorant. On each compact positive-time cylinder
$t\geq t_{\min}>0$, this gives spatial semiconcavity with the
single constant $1/t_{\min}$, uniformly in space and time. Thus,
once the subsolution property below is proved, the hypotheses of
Lemma~\ref{lem:time-lipschitz} and the contact theorem apply locally.

We must check the subsolution property for the changed PDE, rather
than import the zero-viscosity one. At an upper test $\phi$ of $W$,
choose a finite attaining configuration with positive weights and
barycenter $x$. The test for its weighted sum is
\[
 \Phi(Y,t)=\phi\left(\sum_j\theta_jy_j,t\right)
       +\frac{\kappa^2t}{2}
       +\frac1{2t}\sum_j\theta_j
                   \left|y_j-\sum_k\theta_ky_k\right|^2.
\]
Use the synchronized-time theorem of sums from
\cite[Sections 2 and 5]{OpenAI}. Divide each weighted closed jet by
its positive weight, apply its individual subsolution inequality,
and multiply the result back by that weight. Test the block matrix
inequality on the common-time vector $((q_j)_j,1)$, with
\[
 q_j=D_x\phi+\frac{y_j-x}{t},
\]
and on each of the $m$ common spatial-translation vectors
$((e_k)_j,0)$. First let the theorem-of-sums error tend to zero,
then let the time synchronization parameter tend to infinity. The
synchronization Hessian vanishes on all these vectors; large
individual time-gradient components do not enter $G_\tau$. The
resulting inequality is
\[
 D^2\Phi[((q_j)_j,1)]
       +\tau\sum_{k=1}^mD^2\Phi[((e_k)_j,0)]\geq0.
\]
Along the first vector, the centered offsets scale by $1+h/t$,
so their variance divided by $2(t+h)$ is affine in $h$. Along
common spatial translations the variance is constant. Their weighted
velocity is $D_x\phi$. Thus the displayed inequality is exactly
$G_\tau(D\phi,D^2\phi)\geq0$. It proves that $W$ is a subsolution,
including at zero horizontal gradient.

Theorem~\ref{thm:contact} now applies to \eqref{eq:extremal-contact}
on the connected positive-time half-space and gives $W=V$ there.
For any fixed $t>0$, subquadratic spatial growth implies that
$V(y,t)-|y|^2/(2t)$ attains a global maximum at some $y_*$. Average
its maximum inequality over any finite configuration with mean
$y_*$. The variance identity yields
\[
 \sum_j\theta_jV(y_j,t)
      -\frac1{2t}\sum_j\theta_j|y_j-y_*|^2\leq V(y_*,t).
\]
The singleton configuration attains equality, so
$\mathcal C_tV(y_*,t)=V(y_*,t)$. But $W=V$ requires the same gap to
be $\kappa^2t/2>0$ at every point, a contradiction. This proves
Theorem~\ref{thm:viscous}.

\begin{thebibliography}{11}
\raggedright
\small
\setlength{\itemsep}{2pt}
\bibitem{OpenAI}
OpenAI, \emph{Uniform interior $C^{1,\alpha}$ estimates for
infinity-harmonic functions}, October 4, 2026, pinned repository
revision \texttt{fd4aeeb2ee4fc729c18d98444fed42fd0529eeeb},
\href{https://github.com/openai/math/tree/fd4aeeb2ee4fc729c18d98444fed42fd0529eeeb/preprints/Uniform-Interior-C1alpha-Estimates-for-Infinity-Harmonic-Functions-October-4-2026}{manuscript and source files}.
\bibitem{KotschwarNi}
B. Kotschwar and L. Ni, \emph{Local gradient estimates of $p$-harmonic
functions, $1/H$-flow, and an entropy formula},
\href{https://arxiv.org/abs/0711.2291v1}{arXiv:0711.2291v1}.
\bibitem{Uhlenbeck}
K. Uhlenbeck, \emph{Regularity for a class of non-linear elliptic
systems}, Acta Math. \textbf{138} (1977), 219--240,
\href{https://doi.org/10.1007/BF02392316}{doi:10.1007/BF02392316}.
\bibitem{Savin}
O. Savin, \emph{Small perturbation solutions for elliptic equations},
Comm. Partial Differential Equations \textbf{32} (2007), 557--578,
\href{https://doi.org/10.1080/03605300500394405}{doi:10.1080/03605300500394405}.
\bibitem{JulinJuutinen}
V. Julin and P. Juutinen, \emph{A new proof for the equivalence of
weak and viscosity solutions for the $p$-Laplace equation},
\href{https://arxiv.org/abs/1104.2197v1}{arXiv:1104.2197v1}.
\bibitem{Crescimbeni}
R. Crescimbeni, L. Forzani and A. Perini, \emph{A new proof of
Harnack's inequality for elliptic partial differential equations in
divergence form}, Electron. J. Differential Equations (2007), no. 38,
\href{https://ejde.math.txstate.edu/Volumes/2007/38/crescimbeni.pdf}{primary article}.
\bibitem{SanchezMorgado}
H. S\'anchez Morgado, \emph{Time-periodic Evans approach to weak KAM
theory}, \href{https://arxiv.org/abs/1712.10275v1}{arXiv:1712.10275v1}.
\bibitem{PengZhangZhou}
F. Peng, Y. R.-Y. Zhang and Y. Zhou, \emph{Regularity from
$p$-harmonic potentials to $\infty$-harmonic potentials in convex rings},
\href{https://arxiv.org/abs/2310.08093v4}{arXiv:2310.08093v4}.
\bibitem{Bungert}
L. Bungert, \emph{The convergence rate of $p$-harmonic to
infinity-harmonic functions},
\href{https://arxiv.org/abs/2302.08462v2}{arXiv:2302.08462v2}.
\bibitem{SobralTeixeiraUrbano}
A. Sobral, E. V. Teixeira and J. M. Urbano,
\emph{Qualitative-to-quantitative regularity for degenerate equations},
\href{https://arxiv.org/abs/2610.13137v1}{arXiv:2610.13137v1}.
\bibitem{LeeLee}
S.-C. Lee and T. Lee, \emph{The $C^{p'}$-regularity conjecture near $p=2$},
\href{https://arxiv.org/abs/2609.09966v1}{arXiv:2609.09966v1}.
\end{thebibliography}

\end{document}
